\documentclass{article}
\usepackage[utf8]{inputenc}
\usepackage{hyperref}
\usepackage{listings}
\usepackage{xcolor}
\usepackage[margin=1in]{geometry}
\usepackage{booktabs}
\usepackage{longtable}
\usepackage{array}
\usepackage{enumitem}
\usepackage{fancyhdr}
\usepackage{titlesec}

% --- Colour palette ------------------------------------------------------------
\definecolor{codegreen}{rgb}{0,0.6,0}
\definecolor{codegray}{rgb}{0.5,0.5,0.5}
\definecolor{codepurple}{rgb}{0.58,0,0.82}
\definecolor{backcolour}{rgb}{0.95,0.95,0.92}
\definecolor{evblue}{RGB}{30,100,200}
\definecolor{cargoamber}{RGB}{200,120,0}
\definecolor{rowgray}{gray}{0.92}

% --- Listing style -------------------------------------------------------------
\lstdefinestyle{mystyle}{
    backgroundcolor=\color{backcolour},
    commentstyle=\color{codegreen},
    keywordstyle=\color{magenta},
    numberstyle=\tiny\color{codegray},
    stringstyle=\color{codepurple},
    basicstyle=\ttfamily\footnotesize,
    breakatwhitespace=false,
    breaklines=true,
    captionpos=b,
    keepspaces=true,
    numbers=left,
    numbersep=5pt,
    showspaces=false,
    showstringspaces=false,
    showtabs=false,
    tabsize=2,
}
\lstdefinelanguage{json}{
    morestring=[b]",
    morecomment=[l]{//},
    sensitive=false,
}
\lstset{style=mystyle}

% --- HTTP method macros --------------------------------------------------------
\newcommand{\GET}{\colorbox{evblue!20}{\texttt{\textbf{\color{evblue}GET}}}}
\newcommand{\POST}{\colorbox{cargoamber!30}{\texttt{\textbf{\color{cargoamber}POST}}}}
\newcommand{\PUT}{\colorbox{cargoamber!20}{\texttt{\textbf{\color{cargoamber}PUT}}}}

% --- Page header ---------------------------------------------------------------
\pagestyle{fancy}
\fancyhf{}
\rhead{SIAM EV Platform --- Cargo API (External)}
\lhead{\leftmark}
\rfoot{Page \thepage}
\renewcommand{\headrulewidth}{0.4pt}

\hypersetup{
    colorlinks=true,
    linkcolor=evblue,
    urlcolor=evblue,
    pdftitle={SIAM EV Platform - Cargo Space Monitoring API Reference (External)},
    pdfauthor={Ruslee Sutthaweekul}
}

% --- Document ------------------------------------------------------------------
\title{%
    \textbf{SIAM EV Platform}\\[6pt]
    \large Cargo Space Monitoring API Reference\\
    \normalsize External Integration Edition
}
\author{Ruslee Sutthaweekul}
\date{Version 1.0 \quad\textbullet\quad \today}

\begin{document}

\maketitle
\tableofcontents
\newpage

% ==============================================================================
\section{Introduction}
% ==============================================================================

This document describes the REST and WebSocket APIs for the \textbf{Cargo Space
Monitoring} sub-system of the SIAM EV Unified Energy \& Logistics Management
Platform. It is a scoped extract of the full platform API reference, prepared
for handover to an external development partner responsible for the Cargo /
Fleet management system. It does not include the OCPP Charging Monitoring or
Energy Management System (EMS) APIs, which remain internal to SIAM EV; see
Section~\ref{sec:integration-notes} for what a partner integrator needs to know
about those boundaries.

The Cargo Space Monitoring sub-system tracks two things:

\begin{itemize}
    \item \textbf{Trucks} --- mobile units tracked by GPS in real time. Each truck
          carries a load sensor (0--100\%) and a camera that photographs the cargo
          bay automatically when the door closes at departure.
    \item \textbf{Stations} --- static named locations (reference points) where
          trucks arrive and depart. A station has only an identifier, a name, and
          a geographic position.
\end{itemize}

\noindent Cargo photos are stored in S3-compatible object storage and referenced
by time-limited pre-signed URLs.

\subsection{Base URL}

All endpoints are served under the API Gateway. Replace \texttt{\{host\}} with the
deployment address (e.g.\ \texttt{api.siam-ev.example.com}).

\begin{lstlisting}[language=bash, caption=Base URL pattern]
https://{host}/api/v1
\end{lstlisting}

\textbf{Cargo base path:} \texttt{/api/v1/cargo}

\subsection{Authentication}

Every request must carry a Bearer token obtained from the Auth Service.

\begin{lstlisting}[language=bash, caption=Authorization header]
Authorization: Bearer <access_token>
\end{lstlisting}

\begin{center}
\begin{tabular}{>{\ttfamily}lll}
\toprule
\textbf{Header} & \textbf{Required} & \textbf{Description}\\
\midrule
Authorization & Yes & \texttt{Bearer <JWT>} issued by \texttt{/auth/token}\\
Content-Type  & Yes (POST/PUT) & \texttt{application/json}\\
X-Request-ID  & No  & Idempotency key; echoed in response headers\\
\bottomrule
\end{tabular}
\end{center}

\subsection{Common Error Schema}

\begin{lstlisting}[language=json, caption=Error response envelope]
{
  "success": false,
  "error": {
    "code": "TRUCK_NOT_FOUND",
    "message": "No truck with id 'truck-bkk-099' exists.",
    "timestamp": "2026-07-02T10:15:30Z"
  }
}
\end{lstlisting}

\begin{center}
\begin{tabular}{lll}
\toprule
\textbf{HTTP Status} & \textbf{Meaning} & \textbf{Common cause}\\
\midrule
200 & OK & Successful read\\
201 & Created & Resource created\\
400 & Bad Request & Missing or invalid fields\\
401 & Unauthorized & Missing / expired token\\
403 & Forbidden & Insufficient role\\
404 & Not Found & Resource does not exist\\
409 & Conflict & Duplicate resource\\
500 & Internal Error & Upstream / database failure\\
\bottomrule
\end{tabular}
\end{center}

\newpage

% ==============================================================================
\section{Data Models}
% ==============================================================================

\subsection{Truck Object}

\begin{lstlisting}[language=json, caption=Truck object schema]
{
  "id":            "truck-bkk-007",
  "plateNumber":   "GKh-1234",
  "name":          "Truck 07",
  "type":          "StraightTruck",  // StraightTruck | Trailer | Refrigerated | Van
  "powertrain":    "Electric",       // Electric | Diesel | Hybrid
  "energyConsumptionKwhPerKm": 1.35, // rated avg consumption; Electric/Hybrid only, else null
  "capacityCbm":   40,
  "status":        "InTransit",      // Idle | InTransit | Loading | Unloading | Offline
  "currentLoadPct": 72,              // 0-100; from load sensor in cargo bay
  "position": {
    "lat":       13.8562,
    "lng":       100.6211,
    "headingDeg": 245,
    "speedKph":   68,
    "timestamp": "2026-07-02T10:04:55Z"
  },
  "currentTripId": "trip-20260702-0041",
  "driverId":      "usr-drv-019",
  "lastSeenAt":    "2026-07-02T10:04:55Z"
}
\end{lstlisting}

\begin{center}
\begin{tabular}{ll}
\toprule
\textbf{status} & \textbf{Meaning}\\
\midrule
\texttt{Idle}       & At a station, engine off, no active trip\\
\texttt{InTransit}  & Moving between stations\\
\texttt{Loading}    & Door open, cargo being loaded\\
\texttt{Unloading}  & Door open, cargo being unloaded\\
\texttt{Offline}    & GPS/telemetry signal lost \textgreater{} 5 min\\
\bottomrule
\end{tabular}
\end{center}

\noindent\textbf{Note on \texttt{powertrain} / \texttt{energyConsumptionKwhPerKm}:} these fields
identify Electric and Hybrid trucks and their rated energy-consumption rate, and are the
inputs the Cargo Carbon Emissions endpoint (Section~\ref{sec:carbon-emissions}) uses to
compute vehicle-use emissions. \texttt{energyConsumptionKwhPerKm} is \texttt{null} for
Diesel trucks.

\subsection{Route Object}

A \emph{route} defines a named sequence of stations a truck is assigned to serve.

\begin{lstlisting}[language=json, caption=Route object schema]
{
  "id":          "route-bkk-R12",
  "name":        "Bangkok South Loop",
  "description": "Lad Krabang Hub -> Samut Prakan -> Bang Na -> Return",
  "waypoints": [
    { "order": 1, "stationId": "site-bkk-001",
      "name": "Lad Krabang Hub",    "lat": 13.7284, "lng": 100.7498,
      "dwellMinutes": 30 },
    { "order": 2, "stationId": "site-bkk-012",
      "name": "Samut Prakan Depot", "lat": 13.5990, "lng": 100.5964,
      "dwellMinutes": 20 },
    { "order": 3, "stationId": "site-bkk-005",
      "name": "Bang Na Post Office","lat": 13.6607, "lng": 100.6050,
      "dwellMinutes": 15 }
  ],
  "estimatedDurationMin": 130,
  "distanceKm":           68.4,
  "activeOnWorkdays":     true,
  "activeOnWeekends":     false
}
\end{lstlisting}

\subsection{Trip Object}

A \emph{trip} is one end-to-end run of a route by a specific truck.
A cargo photo is captured automatically on each \texttt{DoorClose} event
(i.e.\ when the door shuts after loading at a station).

\begin{lstlisting}[language=json, caption=Trip object schema]
{
  "id":           "trip-20260702-0041",
  "truckId":      "truck-bkk-007",
  "routeId":      "route-bkk-R12",
  "driverId":     "usr-drv-019",
  "status":       "InProgress",  // Scheduled | InProgress | Completed | Cancelled
  "isWorkday":    true,
  "scheduledDepartureAt": "2026-07-02T08:00:00Z",
  "actualDepartureAt":    "2026-07-02T08:07:00Z",
  "actualArrivalAt":      null,  // null while InProgress
  "distanceTraveledKm":   null,  // actual GPS-tracked distance; finalized when status = Completed
  "energyConsumedKwh":    null,  // actual traction energy used; Electric/Hybrid trucks only, else null
  "events": [
    {
      "eventId":    "evt-tr-0041-001",
      "type":       "DoorClose",   // Arrival | Departure | DoorOpen | DoorClose
      "stationId":  "site-bkk-001",
      "stationName":"Lad Krabang Hub",
      "timestamp":  "2026-07-02T08:05:00Z",
      "loadPct":    85,
      "cargoPhoto": {
        "photoId":     "photo-0041-001",
        "url":         "https://storage.siam-ev.example.com/...",
        "urlExpiresAt":"2026-07-02T11:05:00Z",  // pre-signed URL TTL
        "takenAt":     "2026-07-02T08:05:03Z",
        "fileSizeKb":  342
      }
    },
    {
      "eventId":   "evt-tr-0041-002",
      "type":      "Arrival",
      "stationId": "site-bkk-012",
      "stationName":"Samut Prakan Depot",
      "timestamp": "2026-07-02T09:44:00Z",
      "loadPct":   61,
      "cargoPhoto": null
    }
  ]
}
\end{lstlisting}

\subsection{Station Object}

A \emph{station} is a named geographic reference point where trucks pick up and
drop off cargo. It is static master data: the Cargo sub-system stores only the
fields a dispatcher or map view needs to identify and locate the site.

\begin{lstlisting}[language=json, caption=Station object schema]
{
  "id":      "site-bkk-001",
  "name":    "Lad Krabang Hub",
  "address": "999 Moo 1, Lad Krabang, Bangkok 10520",
  "lat":     13.7284,
  "lng":     100.7498,
  "type":    "Hub"                    // Hub | Depot | PostOffice | CustomerSite
}
\end{lstlisting}

\newpage

% ==============================================================================
\section{Truck Endpoints}
% ==============================================================================

\subsection{\GET\ \ \texttt{/cargo/trucks} --- Live truck map data}

Returns all trucks with current GPS position and load percentage for the
real-time map dashboard.

\textbf{Query parameters:}

\begin{center}
\begin{tabular}{>{\ttfamily}llll}
\toprule
\textbf{Param} & \textbf{Type} & \textbf{Default} & \textbf{Description}\\
\midrule
status    & string  & ---   & Filter by status enum\\
routeId   & string  & ---   & Filter trucks assigned to a route\\
limit     & integer & 200   & Max records\\
offset    & integer & 0     & Pagination offset\\
\bottomrule
\end{tabular}
\end{center}

\textbf{Response 200:}
\begin{lstlisting}[language=json, caption=Live truck list response]
{
  "success": true,
  "data": [ /* array of Truck objects */ ],
  "meta": { "total": 18, "inTransit": 11, "idle": 5, "offline": 2 }
}
\end{lstlisting}

\subsection{\GET\ \ \texttt{/cargo/trucks/\{id\}} --- Get single truck}

Returns the full Truck object including the current trip reference.

\textbf{Response 200:}
\begin{lstlisting}[language=json]
{ "success": true, "data": { /* Truck object */ } }
\end{lstlisting}

\subsection{\GET\ \ \texttt{/cargo/trucks/\{id\}/trips} --- Truck trip history}

Returns completed and in-progress trips for a specific truck, with optional
date and day-type filters.

\textbf{Query parameters:}

\begin{center}
\begin{tabular}{>{\ttfamily}llll}
\toprule
\textbf{Param} & \textbf{Type} & \textbf{Required} & \textbf{Description}\\
\midrule
from    & ISO-8601 & No  & Filter trips departing after this time\\
to      & ISO-8601 & No  & Filter trips departing before this time\\
dayType & string   & No  & \texttt{workday} \textbar\ \texttt{weekend} \textbar\ \texttt{all} (default \texttt{all})\\
routeId & string   & No  & Limit to a specific route\\
status  & string   & No  & \texttt{Completed} \textbar\ \texttt{InProgress} \textbar\ \texttt{Cancelled}\\
limit   & integer  & No  & Max records (default 50)\\
offset  & integer  & No  & Pagination offset\\
\bottomrule
\end{tabular}
\end{center}

\textbf{Response 200:}
\begin{lstlisting}[language=json, caption=Truck trip history response]
{
  "success": true,
  "data": [
    {
      "id":                "trip-20260702-0041",
      "routeId":           "route-bkk-R12",
      "routeName":         "Bangkok South Loop",
      "status":            "InProgress",
      "isWorkday":         true,
      "actualDepartureAt": "2026-07-02T08:07:00Z",
      "actualArrivalAt":   null,
      "loadAtDeparturePct": 85,
      "loadAtArrivalPct":   null,
      "eventCount":        2
    }
  ],
  "meta": { "total": 143, "limit": 50, "offset": 0 }
}
\end{lstlisting}

\newpage

% ==============================================================================
\section{Route Endpoints}
% ==============================================================================

\subsection{\GET\ \ \texttt{/cargo/routes} --- List routes}

\textbf{Query parameters:} \texttt{activeOnWorkdays} (bool), \texttt{activeOnWeekends} (bool).

\textbf{Response 200:}
\begin{lstlisting}[language=json]
{ "success": true, "data": [ /* array of Route objects */ ] }
\end{lstlisting}

\subsection{\GET\ \ \texttt{/cargo/routes/\{id\}} --- Route details}

Returns the full Route object including all waypoints with station coordinates.

\textbf{Response 200:}
\begin{lstlisting}[language=json]
{ "success": true, "data": { /* Route object */ } }
\end{lstlisting}

\subsection{\GET\ \ \texttt{/cargo/routes/\{id\}/trips} --- Trips for a route}

Returns all trips that ran (or are running) on a given route, across any truck.
Supports the same \texttt{from}/\texttt{to}/\texttt{dayType} filters as the
truck-level history endpoint.

\textbf{Response 200:}
\begin{lstlisting}[language=json, caption=Route trip history response]
{
  "success": true,
  "data": [
    {
      "id":                "trip-20260702-0041",
      "truckId":           "truck-bkk-007",
      "truckPlate":        "gk-1234",
      "driverId":          "usr-drv-019",
      "status":            "InProgress",
      "isWorkday":         true,
      "actualDepartureAt": "2026-07-02T08:07:00Z",
      "loadAtDeparturePct": 85
    }
  ],
  "meta": { "total": 92, "limit": 50, "offset": 0 }
}
\end{lstlisting}

\newpage

% ==============================================================================
\section{Trip Endpoints}
% ==============================================================================

\subsection{\GET\ \ \texttt{/cargo/trips} --- Search trips}

Global trip search across all trucks and routes with rich filtering --- the
primary endpoint for the history view and the planning dashboard.

\textbf{Query parameters:}

\begin{center}
\begin{tabular}{>{\ttfamily}llll}
\toprule
\textbf{Param} & \textbf{Type} & \textbf{Required} & \textbf{Description}\\
\midrule
from         & ISO-8601 & No  & Departure time lower bound\\
to           & ISO-8601 & No  & Departure time upper bound\\
dayType      & string   & No  & \texttt{workday} \textbar\ \texttt{weekend} \textbar\ \texttt{all}\\
truckId      & string   & No  & Scope to one truck\\
routeId      & string   & No  & Scope to one route\\
stationId    & string   & No  & Trips that visited this station\\
status       & string   & No  & Trip status filter\\
minLoadPct   & integer  & No  & Return trips where departure load $\geq$ value\\
maxLoadPct   & integer  & No  & Return trips where departure load $\leq$ value\\
limit        & integer  & No  & Max records (default 50, max 500)\\
offset       & integer  & No  & Pagination offset\\
\bottomrule
\end{tabular}
\end{center}

\textbf{Response 200:}
\begin{lstlisting}[language=json, caption=Trip search response (summary list)]
{
  "success": true,
  "data": [ /* array of Trip summary objects (no events array) */ ],
  "meta": {
    "total":          320,
    "limit":          50,
    "offset":         0,
    "avgLoadPct":     71.4,
    "workdayCount":   234,
    "weekendCount":   86
  }
}
\end{lstlisting}

\subsection{\GET\ \ \texttt{/cargo/trips/\{id\}} --- Get full trip detail}

Returns the complete Trip object including every event, load reading, and
cargo photo URL for that trip.

\textbf{Response 200:}
\begin{lstlisting}[language=json]
{ "success": true, "data": { /* full Trip object with events array */ } }
\end{lstlisting}

\noindent\textbf{Note on cargo photo URLs:} Pre-signed URLs embedded in the response
expire after 3 hours. Clients should not persist the URL; re-fetch the trip to
get a fresh URL if needed.

\newpage

% ==============================================================================
\section{Planning Analytics Endpoints}
% ==============================================================================

These endpoints aggregate historical trip data into the metrics needed to plan
future route assignments (truck count, trip frequency, load efficiency).

\subsection{\GET\ \ \texttt{/cargo/analytics/utilization} --- Load utilization over time}

Returns average load percentage and trip frequency bucketed by time interval
and optionally split by route, truck, or day type.

\textbf{Query parameters:}

\begin{center}
\begin{tabular}{>{\ttfamily}llll}
\toprule
\textbf{Param} & \textbf{Type} & \textbf{Required} & \textbf{Description}\\
\midrule
from      & ISO-8601 & Yes & Start of analysis window\\
to        & ISO-8601 & Yes & End of analysis window\\
groupBy   & string   & No  & \texttt{route} \textbar\ \texttt{truck} \textbar\ \texttt{dayType} \textbar\ \texttt{hour}\\
routeId   & string   & No  & Scope to one route\\
dayType   & string   & No  & \texttt{workday} \textbar\ \texttt{weekend} \textbar\ \texttt{all}\\
interval  & string   & No  & Bucket size: \texttt{1d} \textbar\ \texttt{1w} \textbar\ \texttt{1mo} (default \texttt{1d})\\
\bottomrule
\end{tabular}
\end{center}

\textbf{Response 200:}
\begin{lstlisting}[language=json, caption=Utilization analytics response]
{
  "success": true,
  "data": {
    "window": { "from": "2026-06-01T00:00:00Z", "to": "2026-07-01T00:00:00Z" },
    "interval": "1d",
    "groupBy":  "route",
    "series": [
      {
        "groupKey":    "route-bkk-R12",
        "groupLabel":  "Bangkok South Loop",
        "buckets": [
          {
            "date":          "2026-06-01",
            "tripsCompleted": 3,
            "avgLoadPct":    74.2,
            "peakLoadPct":   95,
            "minLoadPct":    42,
            "totalKm":       205.2
          }
        ]
      }
    ]
  }
}
\end{lstlisting}

\subsection{\GET\ \ \texttt{/cargo/analytics/planning} --- Fleet planning recommendations}

Analyses historical patterns and returns data-driven recommendations for
route scheduling, truck count, and travel frequency. This is the primary
output for the planning workflow.

\textbf{Query parameters:}

\begin{center}
\begin{tabular}{>{\ttfamily}llll}
\toprule
\textbf{Param} & \textbf{Type} & \textbf{Required} & \textbf{Description}\\
\midrule
from       & ISO-8601 & Yes & Historical window start\\
to         & ISO-8601 & Yes & Historical window end\\
routeId    & string   & No  & Limit analysis to one route\\
targetUtil & number   & No  & Desired avg load \% (default 80)\\
\bottomrule
\end{tabular}
\end{center}

\textbf{Response 200:}
\begin{lstlisting}[language=json, caption=Planning recommendations response]
{
  "success": true,
  "data": {
    "analysisWindow": { "from": "2026-06-01T00:00:00Z", "to": "2026-07-01T00:00:00Z" },
    "targetUtilPct":  80,
    "routes": [
      {
        "routeId":    "route-bkk-R12",
        "routeName":  "Bangkok South Loop",
        "workdays": {
          "avgTripsPerDay":   3.1,
          "avgLoadPct":       74.2,
          "peakDemandHour":   "08:00",
          "currentTrucksUsed": 2,
          "recommendedTrucks": 2,
          "recommendedFreqPerDay": 3,
          "utilizationGapPct": 5.8
        },
        "weekends": {
          "avgTripsPerDay":   1.0,
          "avgLoadPct":       41.3,
          "peakDemandHour":   "10:00",
          "currentTrucksUsed": 2,
          "recommendedTrucks": 1,
          "recommendedFreqPerDay": 1,
          "utilizationGapPct": 38.7
        },
        "insights": [
          "Weekend demand is 55% lower than workday demand.",
          "Reallocating 1 truck from weekends could serve 2 additional workday trips.",
          "Peak load occurs between 07:00 and 09:00 on workdays."
        ]
      }
    ]
  }
}
\end{lstlisting}

\subsection{\GET\ \ \texttt{/cargo/analytics/load-history} --- Load trend with photos}

Returns load readings paired with cargo photo thumbnails for a route or truck
over a date range. Designed for the history gallery view in the dashboard.

\textbf{Query parameters:}

\begin{center}
\begin{tabular}{>{\ttfamily}llll}
\toprule
\textbf{Param} & \textbf{Type} & \textbf{Required} & \textbf{Description}\\
\midrule
from     & ISO-8601 & Yes & Start date\\
to       & ISO-8601 & Yes & End date\\
routeId  & string   & No  & Filter by route\\
truckId  & string   & No  & Filter by truck\\
dayType  & string   & No  & \texttt{workday} \textbar\ \texttt{weekend} \textbar\ \texttt{all}\\
limit    & integer  & No  & Max events returned (default 100)\\
\bottomrule
\end{tabular}
\end{center}

\textbf{Response 200:}
\begin{lstlisting}[language=json, caption=Load history with photos response]
{
  "success": true,
  "data": [
    {
      "tripId":       "trip-20260702-0041",
      "truckId":      "truck-bkk-007",
      "truckPlate":   "gk-1234",
      "routeName":    "Bangkok South Loop",
      "stationName":  "Lad Krabang Hub",
      "departureAt":  "2026-07-02T08:07:00Z",
      "isWorkday":    true,
      "loadPct":      85,
      "photo": {
        "photoId":      "photo-0041-001",
        "url":          "https://storage.siam-ev.example.com/...",
        "thumbnailUrl": "https://storage.siam-ev.example.com/.../thumb",
        "urlExpiresAt": "2026-07-02T11:07:00Z",
        "takenAt":      "2026-07-02T08:05:03Z"
      }
    }
  ],
  "meta": { "total": 187, "limit": 100, "offset": 0 }
}
\end{lstlisting}

\newpage

% ==============================================================================
\section{Carbon Emissions \& Reduction Endpoints}
\label{sec:carbon-emissions}
% ==============================================================================

The Cargo sub-system owns the raw driving data --- distance and, for Electric/Hybrid
trucks, energy consumption (Section~2.1) --- needed to calculate vehicle-use carbon
emissions and the emission reduction achieved versus an equivalent internal-combustion
(ICE) truck. It does \emph{not} own the emission-factor reference values themselves
(grid carbon intensity, ICE baseline rate); those are owned by SIAM EV's internal
Carbon Credit / EMS system and must be supplied by the caller as query parameters on
each request. Charging-session emissions and building/facility emissions are computed
entirely within the EMS and are not exposed via the Cargo API --- see
Section~\ref{sec:integration-notes}.

\subsection{\GET\ \ \texttt{/cargo/analytics/carbon-emissions} --- Vehicle-use emissions and reduction}

Aggregates distance and energy consumption for Electric/Hybrid trucks over a time
window and applies caller-supplied emission factors to compute emissions and the
reduction versus an ICE baseline.

\textbf{Query parameters:}

\begin{center}
\begin{tabular}{>{\ttfamily}llll}
\toprule
\textbf{Param} & \textbf{Type} & \textbf{Required} & \textbf{Description}\\
\midrule
from   & ISO-8601 & Yes & Start of analysis window\\
to     & ISO-8601 & Yes & End of analysis window\\
gridEmissionFactorKgPerKwh & number & Yes & Electricity carbon intensity (kgCO\textsubscript{2}e/kWh) applied to truck energy consumption\\
baselineEmissionFactorKgPerKm & number & Yes & Reference emission rate (kgCO\textsubscript{2}e/km) of an equivalent ICE truck, used for the reduction calc\\
routeId   & string & No & Scope to one route\\
truckId   & string & No & Scope to one truck\\
groupBy   & string & No & \texttt{route} \textbar\ \texttt{truck} \textbar\ \texttt{dayType} (default \texttt{route})\\
dayType   & string & No & \texttt{workday} \textbar\ \texttt{weekend} \textbar\ \texttt{all}\\
\bottomrule
\end{tabular}
\end{center}

\textbf{Response 200:}
\begin{lstlisting}[language=json, caption=Carbon emissions and reduction response]
{
  "success": true,
  "data": {
    "window": { "from": "2026-06-01T00:00:00Z", "to": "2026-07-01T00:00:00Z" },
    "groupBy": "route",
    "gridEmissionFactorKgPerKwh":     0.4999,
    "baselineEmissionFactorKgPerKm":  0.85,
    "series": [
      {
        "groupKey":                "route-bkk-R12",
        "groupLabel":              "Bangkok South Loop",
        "tripsIncluded":           88,
        "distanceKm":              4108.4,
        "energyConsumedKwh":       5546.3,
        "evEmissionsKgCO2e":       2772.9,
        "baselineEmissionsKgCO2e": 3492.1,
        "emissionReductionKgCO2e": 719.2
      }
    ],
    "totals": {
      "distanceKm":              4108.4,
      "energyConsumedKwh":       5546.3,
      "evEmissionsKgCO2e":       2772.9,
      "baselineEmissionsKgCO2e": 3492.1,
      "emissionReductionKgCO2e": 719.2
    }
  }
}
\end{lstlisting}

\noindent\textbf{Calculation:}
\begin{itemize}[leftmargin=*]
    \item \texttt{evEmissionsKgCO2e} = $\sum$ \texttt{energyConsumedKwh} (actual, per
          completed trip) $\times$ \texttt{gridEmissionFactorKgPerKwh}; if a trip's
          actual \texttt{energyConsumedKwh} is null, it falls back to
          \texttt{distanceTraveledKm} $\times$ truck \texttt{energyConsumptionKwhPerKm}.
    \item \texttt{baselineEmissionsKgCO2e} = \texttt{distanceKm} $\times$
          \texttt{baselineEmissionFactorKgPerKm}.
    \item \texttt{emissionReductionKgCO2e} = \texttt{baselineEmissionsKgCO2e} $-$
          \texttt{evEmissionsKgCO2e}.
\end{itemize}

\noindent Only trips completed by Electric/Hybrid trucks (\texttt{powertrain} $\neq$
\texttt{Diesel}) are included in \texttt{evEmissionsKgCO2e}; \texttt{distanceKm} in the
response includes all trips in scope regardless of powertrain, so the baseline figure
represents the counterfactual of running the same distance with ICE trucks.

\newpage

% ==============================================================================
\section{Station Endpoints}
% ==============================================================================

Stations are static master data. These endpoints populate map markers, filter
dropdowns, and the route builder. They carry no real-time energy or
environmental data.

\subsection{\GET\ \ \texttt{/cargo/stations} --- List stations}

\textbf{Query parameters:} \texttt{type}, \texttt{limit}, \texttt{offset}.

\textbf{Response 200:}
\begin{lstlisting}[language=json]
{
  "success": true,
  "data": [ /* array of Station objects */ ],
  "meta": { "total": 18 }
}
\end{lstlisting}

\subsection{\GET\ \ \texttt{/cargo/stations/\{id\}} --- Get station}

Returns the Station object (id, name, address, lat, lng, type).

\section{Alert Endpoints}

\subsection{\GET\ \ \texttt{/cargo/alerts} --- List active alerts}

\textbf{Query parameters:} \texttt{stationId}, \texttt{truckId}, \texttt{severity}
(\texttt{Critical} \textbar\ \texttt{Warning}), \texttt{limit}, \texttt{offset}.

\textbf{Response 200:}
\begin{lstlisting}[language=json, caption=Active alerts response]
{
  "success": true,
  "data": [
    {
      "alertId":       "alrt-20260702-007",
      "sourceType":    "Zone",          // Zone | Truck
      "sourceId":      "zone-stn-001-cold-01",
      "type":          "TemperatureExceeded",
      "severity":      "Critical",
      "value":         9.1,
      "threshold":     8.0,
      "unit":          "degC",
      "triggeredAt":   "2026-07-02T10:22:00Z",
      "acknowledgedAt": null
    }
  ]
}
\end{lstlisting}

\noindent\textbf{Note on \texttt{sourceType}:} \texttt{Truck} alerts originate from
sensors on the vehicle (e.g.\ a refrigerated compartment). \texttt{Zone} alerts
originate from station-level environmental monitoring that is configured and
owned by an internal SIAM EV system outside this document's scope --- treat
\texttt{sourceId} as an opaque string for display/filtering purposes.

\subsection{\POST\ \ \texttt{/cargo/alerts/\{alertId\}/acknowledge} --- Acknowledge alert}

\textbf{Request body:}
\begin{lstlisting}[language=json]
{ "acknowledgedBy": "usr-0007", "note": "Compressor maintenance dispatched." }
\end{lstlisting}

\newpage

% ==============================================================================
\section{WebSocket --- Live Cargo and Fleet Feed}
% ==============================================================================

A single WebSocket connection delivers both fleet (truck) and station (zone)
events. Use the \texttt{scope} parameter to subscribe only to the events your
dashboard panel needs.

\begin{lstlisting}[language=bash]
wss://{host}/api/v1/cargo/ws?scope=trucks,zones&token=<access_token>
\end{lstlisting}

\textbf{Truck position and load update:}
\begin{lstlisting}[language=json, caption=TruckPositionUpdate event]
{
  "event":   "TruckPositionUpdate",
  "truckId": "truck-bkk-007",
  "data": {
    "lat":        13.8562,
    "lng":        100.6211,
    "headingDeg": 245,
    "speedKph":   68,
    "loadPct":    72,
    "timestamp":  "2026-07-02T10:04:55Z"
  }
}
\end{lstlisting}

\textbf{Cargo photo ready event} (emitted seconds after door closes at a station):
\begin{lstlisting}[language=json, caption=CargoPhotoReady event]
{
  "event":   "CargoPhotoReady",
  "tripId":  "trip-20260702-0041",
  "truckId": "truck-bkk-007",
  "data": {
    "eventId":      "evt-tr-0041-001",
    "stationId":    "site-bkk-001",
    "stationName":  "Lad Krabang Hub",
    "loadPct":      85,
    "photoId":      "photo-0041-001",
    "thumbnailUrl": "https://storage.siam-ev.example.com/.../thumb",
    "takenAt":      "2026-07-02T08:05:03Z"
  }
}
\end{lstlisting}

\textbf{All WebSocket event types:}

\begin{center}
\begin{tabular}{lll}
\toprule
\textbf{event} & \textbf{scope} & \textbf{Trigger}\\
\midrule
\texttt{TruckPositionUpdate} & trucks & GPS reading every 5 s while InTransit\\
\texttt{TruckStatusChange}   & trucks & Status transitions (e.g.\ InTransit $\to$ Idle)\\
\texttt{TripStarted}         & trucks & New trip begins (truck departs station)\\
\texttt{TripCompleted}       & trucks & Trip ends at final waypoint\\
\texttt{CargoPhotoReady}     & trucks & Photo captured on DoorClose; thumbnail available\\
\texttt{EnvironmentUpdate}   & zones  & Zone sensor reading every 30 s\\
\texttt{AlertTriggered}      & zones  & Threshold breach on zone or truck sensor\\
\texttt{AlertResolved}       & zones  & Reading returns within threshold\\
\bottomrule
\end{tabular}
\end{center}

\newpage

% ==============================================================================
\section{Integration Notes}
\label{sec:integration-notes}
% ==============================================================================

This section summarizes, at a high level, how the Cargo system fits into the
wider SIAM EV platform. Details of the other sub-systems are intentionally
excluded from this handover document.

\begin{itemize}[leftmargin=*]
    \item \textbf{Station identity is shared internally.} A \texttt{stationId}
          (e.g.\ \texttt{site-bkk-001}) may correspond to the same physical site
          in other internal SIAM EV systems. This cross-reference is coordinated
          entirely on the SIAM EV side; the Cargo API contract in this document
          does not change based on it, and no additional integration work is
          required by the Cargo system to support it.
    \item \textbf{Zone-sourced alerts.} As noted in Section~6, some entries returned
          by \texttt{GET /cargo/alerts} originate from station environmental
          monitoring that is owned and configured internally. The Cargo system
          only needs to display and allow acknowledgement of these alerts via the
          documented endpoints --- it does not need to compute or validate them.
    \item \textbf{Carbon accounting is split by data ownership.} The Cargo system
          computes vehicle-use emissions and emission reduction
          (Section~\ref{sec:carbon-emissions}) because it owns the underlying
          telemetry (distance, truck energy consumption) --- but it does not own the
          emission-factor reference values, which the caller must supply. Charging-session
          emissions and building/facility emissions are calculated entirely inside the
          internal EMS and are not exposed via any Cargo endpoint.
    \item \textbf{Demand-response and energy coordination.} SIAM EV runs internal
          workflows that may temporarily prioritize certain stations (e.g.\ during
          a facility power event). These do not require any Cargo API changes or
          calls to systems outside this document.
    \item \textbf{Out of scope for this handover.} OCPP charging telemetry, the
          Energy Management System (EMS) API, and the outbound webhook registry
          are separate SIAM EV--internal systems and are not included here. If a
          future feature requires the Cargo system to call one of those APIs
          directly, SIAM EV will provide a supplementary interface specification
          at that time.
\end{itemize}

\newpage

% ==============================================================================
\section{API Quick Reference}
% ==============================================================================

\begin{longtable}{>{\ttfamily\footnotesize}p{0.45\textwidth} p{0.08\textwidth} p{0.38\textwidth}}
\toprule
\textbf{Endpoint} & \textbf{Method} & \textbf{Description}\\
\midrule
\endfirsthead
\toprule
\textbf{Endpoint} & \textbf{Method} & \textbf{Description}\\
\midrule
\endhead
\bottomrule
\endfoot

\multicolumn{3}{l}{\textit{Fleet}}\\[2pt]
/cargo/trucks                            & GET & Live truck positions and loads (map)\\
/cargo/trucks/\{id\}                     & GET & Single truck detail\\
/cargo/trucks/\{id\}/trips               & GET & Trip history for a truck\\
/cargo/routes                            & GET & List routes\\
/cargo/routes/\{id\}                     & GET & Route detail with waypoints\\
/cargo/routes/\{id\}/trips               & GET & All trips on a route\\
/cargo/trips                             & GET & Global trip search with date/dayType filter\\
/cargo/trips/\{id\}                      & GET & Full trip detail with events and photos\\[4pt]

\multicolumn{3}{l}{\textit{Analytics}}\\[2pt]
/cargo/analytics/utilization             & GET & Load \% and trip frequency over time\\
/cargo/analytics/planning                & GET & Fleet sizing and frequency recommendations\\
/cargo/analytics/load-history            & GET & Load trend with cargo photo gallery\\
/cargo/analytics/carbon-emissions        & GET & Vehicle-use carbon emissions and ICE-baseline reduction\\[4pt]

\multicolumn{3}{l}{\textit{Stations (static reference data)}}\\[2pt]
/cargo/stations                          & GET & List stations (name, position, type)\\
/cargo/stations/\{id\}                   & GET & Single station\\
/cargo/alerts                            & GET & Active alerts (trucks and stations)\\
/cargo/alerts/\{alertId\}/acknowledge    & POST & Acknowledge an alert\\
/cargo/ws                                & WS  & Live truck and station event feed\\
\end{longtable}

\end{document}
